% ===================================================================
% Appendix: Analytic study of the temporal-sheaf core (synthetic + per-event)
% Numbers: results/analytic/synth_v4/{synth_*/meta.json, runs/*/summary.json},
%          results/analytic/forum_factorial/*/summary.json, results/analytic/DESIGN.md
% ===================================================================
\section{What the Temporal-Sheaf Core Captures: A Controlled Study}
\label{app:analytic}

The benchmark ablations (App.~\ref{app:faithful:results}) attribute
performance to components under real data, where a recurrency head can
absorb most of the temporal signal. This appendix asks the complementary
question under controlled conditions: which aspects of a temporal graph
\emph{can} the temporal-sheaf core capture, and does it capture them in
practice? We use a synthetic generator with known ceilings, then return to
real data with a per-event analysis.

\paragraph{Generator.}
$n{=}100$ nodes in $k{=}4$ communities carry unit-norm latent states
$x_u(t)\in\mathbb{R}^4$; each community pair $(a,b)$ has a planted
orthogonal transport $O_{ab}$ (the sheaf). States rotate at
community-specific rates (0.01--0.05 rad per time unit; ``static'' regime:
zero rates) and receive a small kick after each event. Events arrive with
bursty (Pareto) gaps; with probability $\rho{=}0.1$ a source repeats a past
partner, otherwise it picks a partner with probability
$\propto\exp(-\|O_{ab}x_u-x_v\|^2/\tau)$, $\tau{=}0.03$. Node attributes
expose $x_u(0)$ and the community, so a model must learn the rates and the
transports from events (system identification) rather than infer hidden
states. 100k events; 70/15/15 split in time; snapshots of 20 time units;
full-vocabulary MRR over all $n{-}1$ candidates.

\paragraph{Oracle ceilings.}
Before training anything we score the generator's own kernels on the test
period, which bounds what each mechanism can contribute:

\begin{table}[h]
\centering\small
\caption{Oracle MRR on the synthetic test period (chance $=0.010$).
``Static features'' knows $x(0)$ and the transports; ``last snapshot''
knows the true state at the previous snapshot boundary but does not
extrapolate the gap; ``event time'' knows the true state at the event.}
\begin{tabular}{lcccc}
\toprule
Regime & no transport & static features & last snapshot & event time \\
\midrule
static (rates 0)          & 0.196 & 0.767 & 0.875 & 0.875 \\
evolving, rotations       & 0.186 & 0.178 & 0.461 & 0.821 \\
evolving, identity maps   & 0.823 & 0.179 & 0.632 & 0.823 \\
\bottomrule
\end{tabular}
\end{table}

The ladder in the evolving regime (0.18 $<$ 0.46 $<$ 0.82) is exactly the
one the core is built to climb: a persistent per-node state (memory)
reaches the second rung, and extrapolating that state across the physical
gap $\Delta_k$ reaches the third. In the static regime the task reduces to
learning the transports. An i.i.d.-trained bilinear scorer over
$x(0)\otimes\text{community}$ features reaches 0.807 of the 0.875 static
ceiling but only 0.200 in the evolving regime, confirming that the evolving
signal is temporal.

\paragraph{Results.}
All variants share the recurrency head, the observed attributes and a
32-channel width, trained 8 epochs under the leak-free protocol (seed 43;
test MRR over all candidates):

\begin{table}[h]
\centering\small
\caption{Synthetic study, test MRR. ``Efficient'' is the released
implementation (no persistent recurrence, no $\Delta_k$).}
\begin{tabular}{lccccc}
\toprule
Regime & full & identity maps & no $\Delta_k$ & current-only maps & efficient \\
\midrule
static                  & 0.849 & 0.846 & 0.849 & 0.847 & 0.790 \\
evolving, rotations     & 0.276 & 0.277 & 0.276 & 0.276 & 0.176 \\
evolving, identity maps & 0.290 & 0.290 & 0.288 & 0.289 & 0.200 \\
\bottomrule
\end{tabular}
\end{table}

Two conclusions. (i) The faithful architecture beats the efficient
implementation by 0.06--0.10 MRR in every regime and reaches 0.85 of the
0.875 static ceiling. (ii) Its variants are indistinguishable (within
0.003): with a recurrency head present, none of the faithful model's own
mechanisms is exercised, and in the evolving regime the model stops between
the ``static features'' rung (0.18--0.20) and the ``last snapshot'' rung
(0.46), i.e.\ it does not extrapolate the planted dynamics across gaps.
Compressing or stretching the test-period gaps by $4\times$ moves the full
model by only $-0.012$ / $+0.005$. The synthetic study is therefore a
bounded negative for the mechanisms and a positive for the architecture:
sheaf diffusion over the current event graph, with the observed features and
the head, carries the gain, and the temporal core's memory and gap
extrapolation do not add to it within this budget.

\paragraph{Back to real data.}
The per-event analysis of Table~\ref{tab:stratified} and the head-vs-core
factorial of Table~\ref{tab:factorial} on \textsc{thgl-forum} give the
positive side of the same question: when the head is removed the core is
worth $+0.15$ MRR and the learned restriction maps are the mechanism, and
with the head present the core's contribution is confined to novel partners
and grows with the inter-event gap. Together, the two studies locate the
temporal-sheaf core's value where recurrence heuristics have nothing to
say: predicting \emph{new} interactions from the geometry of the current
event graph.
